\documentclass[11pt,a4paper]{article}
\usepackage{iftex}
\ifPDFTeX
  \usepackage[T1]{fontenc}
  \usepackage[utf8]{inputenc}
  \usepackage{lmodern}
  \input{glyphtounicode}
  \pdfgentounicode=1
\else
  \usepackage{fontspec}
  % Kerning off: kerned pairs ("A WS", "T ree") split words for ATS text extraction.
  \setmainfont{lmroman10-regular}[Extension=.otf, BoldFont=lmroman10-bold,
    ItalicFont=lmroman10-italic, BoldItalicFont=lmroman10-bolditalic, Kerning=Off]
\fi
\usepackage[margin=0.9in]{geometry}
\usepackage[hidelinks]{hyperref}
\hypersetup{pdftitle={Abhinav Kaushik - Cover Letter - Sonar}, pdfauthor={Abhinav Kaushik}}

\pagestyle{empty}
\setlength{\parindent}{0pt}
\setlength{\parskip}{10pt}
\raggedright

\begin{document}

{\Large\bfseries Abhinav Kaushik}\\
AI Research Engineer\\
New Delhi, India \textbar{} +91 8700571859 \textbar{} \href{mailto:aabhinav.kaushik@gmail.com}{aabhinav.kaushik@gmail.com} \textbar{} \href{https://www.linkedin.com/in/abhinav-kaushik10/}{LinkedIn}  

September 27, 2026

Hiring Team\\
Sonar

\textbf{Re: AI Research Engineer}

Dear Hiring Team,

I am applying for the AI Research Engineer role in the Data \& Agentic team. My work sits between research and product: I build LLM and RAG systems, fine-tune smaller models where they help, and set up the evaluations that tell us whether an idea is worth shipping. Most of it started as a quick prototype and ended up in production.

Code verification matters to me because more and more code is written by AI assistants, and that only works when the review behind it can be trusted. There is also a personal reason: my fiancée lives in Berlin, and we plan to move to Geneva, Switzerland together and settle there long term.

I would bring a practical research habit. I fine-tuned a cross-encoder on positive and negative pairs to improve retrieval, and built evaluation pipelines with LLM evaluators and tracing that cut human review by 43\%. At a previous role I prototyped 15+ AI applications, and 7 of them moved into production or paid client work.

Alongside work I am building ThinkTree.AI, a non-profit AI learning platform used by NGO teachers to support students with ADHD. It reminds me that research only matters once it works for real users. I would be happy to talk about how I could contribute.

Sincerely,\\[4pt]
Abhinav Kaushik

\end{document}